% source: Robin Hartshorne, "Algebraic Geometry", GTM 52 (1977)
% pdf_page: 0155
% doc_page: 138 (Chapter II, Section 6, "Divisors")
% retrieved: 2026-07-17
% transcribed_by: claude-fable-5 (vision, rendered at 200dpi via pdftoppm)

% running head: II Schemes / p. 138

\DeclareMathOperator{\Spec}{Spec}
\DeclareMathOperator{\Cl}{Cl}

\begin{proposition}[6.9]
Let $f\colon X \to Y$ be a finite morphism of nonsingular curves. Then for
any divisor $D$ on $Y$ we have $\deg f^*D = \deg f \cdot \deg D$.
\end{proposition}

\begin{proof}
It will be sufficient to show that for any closed point $Q \in Y$ we have
$\deg f^*Q = \deg f$. Let $V = \Spec B$ be an open affine subset of $Y$
containing $Q$. Let $A$ be the integral closure of $B$ in $K(X)$. Then, as
in the proof of (6.8), $U = \Spec A$ is the open subset $f^{-1}V$ of $X$.
Let $\mathfrak{m}_Q$ be the maximal ideal of $Q$ in $B$. We localize both
$B$ and $A$ with respect to the multiplicative system
$S = B - \mathfrak{m}_Q$, and we obtain a ring extension
$\mathcal{O}_Q \hookrightarrow A'$, where $A'$ is a finitely generated
$\mathcal{O}_Q$-module. Now $A'$ is torsion-free, and has rank equal to
$r = [K(X)\colon K(Y)]$, so $A'$ is a free $\mathcal{O}_Q$-module of rank
$r = \deg f$. If $t$ is a local parameter at $Q$, it follows that $A'/tA'$
is a $k$-vector space of dimension $r$.

On the other hand, the points $P_i$ of $X$ such that $f(P_i) = Q$ are in
1-1 correspondence with the maximal ideals $\mathfrak{m}_i$ of $A'$, and for
each $i$, $A'_{\mathfrak{m}_i} = \mathcal{O}_{P_i}$. Clearly
$tA' = \bigcap_i (tA'_{\mathfrak{m}_i} \cap A')$, so by the Chinese
remainder theorem,
\[
  \dim_k A'/tA' = \sum_i \dim_k A'/(tA'_{\mathfrak{m}_i} \cap A').
\]
But
\[
  A'/(tA'_{\mathfrak{m}_i} \cap A') \cong
  A'_{\mathfrak{m}_i}/tA'_{\mathfrak{m}_i} =
  \mathcal{O}_{P_i}/t\mathcal{O}_{P_i},
\]
so the dimensions in the sum above are just equal to $v_{P_i}(t)$. But
$f^*Q = \sum v_{P_i}(t) \cdot P_i$, so we have shown that
$\deg f^*Q = \deg f$ as required.
\end{proof}

\begin{corollary}[6.10]
A principal divisor on a complete nonsingular curve $X$ has degree zero.
Consequently the degree function induces a surjective homomorphism
$\deg\colon \Cl X \to \mathbf{Z}$.
\end{corollary}

\begin{proof}
Let $f \in K(X)^*$. If $f \in k$, then $(f) = 0$, so there is nothing to
prove. If $f \notin k$, then the inclusion of fields
$k(f) \subseteq K(X)$ induces a finite morphism
$\varphi\colon X \to \mathbf{P}^1$. It is a morphism by (I, 6.12), and it is
finite by (6.8). Now $(f) = \varphi^*(\{0\} - \{\infty\})$. Since
$\{0\} - \{\infty\}$ is a divisor of degree 0 on $\mathbf{P}^1$, we conclude
that $(f)$ has degree 0 on $X$.

Thus the degree of a divisor on $X$ depends only on its linear equivalence
class, and we obtain a homomorphism $\Cl X \to \mathbf{Z}$ as stated. It is
surjective, because the degree of a single point is 1.
\end{proof}

\begin{example}[6.10.1]
A complete nonsingular curve $X$ is rational if and only if there exist two
distinct points $P,Q \in X$ with $P \sim Q$. Recall that \emph{rational}
means birational to $\mathbf{P}^1$. If $X$ is rational, then in fact it is
isomorphic to $\mathbf{P}^1$ by (6.7). And on $\mathbf{P}^1$ we have already
seen that any two points are linearly equivalent (6.4). Conversely, suppose
$X$ has two points $P \neq Q$ with $P \sim Q$. Then there is a rational
function $f \in K(X)$ with $(f) = P - Q$. Consider the morphism
$\varphi\colon X \to \mathbf{P}^1$ determined by $f$ as in the proof of
(6.10). We have $\varphi^*(\{0\}) = P$, so $\varphi$ must be a morphism of
degree 1. In other words, $\varphi$ is birational, so $X$ is rational.
\end{example}
